% Formula sheet, Topics 1-2B (Prof. Ricca, Finanças II). Full version.
% Built from Notas/L01-07, L02-08, L03-06 (formula notes) plus the TVM/bond formulas used in lists
% R2, R2b and PS1-PS2. Compile: pdflatex twice; check pdffonts (Type 1, uni = yes).
\documentclass[10pt,a4paper]{article}
\usepackage[utf8]{inputenc}
\usepackage[T1]{fontenc}
\usepackage{lmodern}   % Type 1 vector fonts (replaces PK bitmaps)
\usepackage{cmap}      % ToUnicode CMap -> copyable text and accents
\usepackage[english]{babel}
\usepackage{amsmath,amssymb}
\usepackage[a4paper,margin=1.8cm]{geometry}
\usepackage{enumitem}
\usepackage{xcolor}
\usepackage{fancyhdr}
\usepackage{microtype}
\setlist{nosep,leftmargin=1.2em}
\pagestyle{fancy}\fancyhf{}
\lhead{\footnotesize Finan\c{c}as II --- Corporativas (MPE Insper)}
\rhead{\footnotesize Formula sheet --- Topics 1 to 2B}
\cfoot{\footnotesize\thepage}
\renewcommand{\headrulewidth}{0.3pt}
\setlength{\headheight}{14pt}
\setlength{\parindent}{0pt}
\setlength{\parskip}{3pt}
\newcommand{\blk}[1]{\subsection*{\color{blue!55!black}#1}}
\newcommand{\idea}{\textbf{Concept.}~}
\newcommand{\warn}{\textbf{Watch out.}~}
\newcommand{\E}{\mathbb{E}}

\title{Formula sheet --- Topics 1, 2A and 2B\\[2pt]
\large Accounting, valuation and the frictionless benchmark, taxes and the trade-off}
\author{Finan\c{c}as II --- Corporativas (MPE Insper)}
\date{}

\begin{document}
\maketitle
\vspace{-1.5em}
\noindent Notation: $A$ assets, $L$ liabilities, $D$ debt, $E$ equity (book unless marked mkt), $V=D+E$,
$\tau$ corporate tax rate, $r_f$ risk-free rate, $r_D$, $r_E$, $r_A$ expected returns on debt,
equity, assets; $r_U$ unlevered cost of capital; $g$ growth; $N$ shares; $X$ operating cash flow.
Block numbers (1.2, 2A.8, \dots) follow the slide blocks of \texttt{Map/lectures-index.md}.

% =====================================================================
\section*{Topic 1 --- Accounting review and the capital structure}

\blk{1.2 Balance-sheet view of the firm}
\begin{gather*}
A = L + E,\qquad \Delta A=\Delta L+\Delta E,\qquad
\underbrace{\min(X,F)}_{\text{creditors}}+\underbrace{\max(X-F,0)}_{\text{owners}} = X,\\
FCF_A = FCF_L + FCF_E \;\Rightarrow\; A^{\text{mkt}} = L^{\text{mkt}} + E^{\text{mkt}}.
\end{gather*}
\idea The identity holds in book \emph{and} market value. Liabilities are senior and specific;
equity is junior and residual. Firms sell goods and services \emph{and promises}. Recognition:
tangibles at historical cost; most intangibles are not on the balance sheet (about 90\% of S\&P 500
EV).

\blk{1.3 Balance-sheet items}
\begin{gather*}
NWC = CA - CL,\qquad \text{Dep (straight line)}=\frac{K_0-S}{n},\qquad
BV_k = K_0-\sum_{j\le k}\text{Dep}_j,\\
\text{Goodwill}=\text{price paid}-\text{fair value of net identifiable assets}.
\end{gather*}
\idea Current = converted to cash / due within one year. Depreciation is \emph{not} a cash expense.
Trade payables are short-term finance (NWC can be negative).

\blk{1.4--1.5 Book vs.\ market, leverage and liquidity}
\begin{gather*}
E^{\text{mkt}} = P\times N,\qquad \frac{M}{B}=\frac{E^{\text{mkt}}}{E^{\text{book}}},\qquad
\text{Net debt}=D-\text{Cash},\qquad EV = E^{\text{mkt}}+D-\text{Cash},\\
\text{Current ratio}=\frac{CA}{CL},\qquad \text{Quick ratio}=\frac{CA-\text{Inv}}{CL}.
\end{gather*}
\idea Market values are forward-looking, book values backward-looking; book equity can be negative
(Starbucks) while market equity is large. Leverage: $D/E$ book vs.\ market. High $M/B$ = growth
stock, low = value stock. Covenants are written on book ratios.

\blk{1.6 Income statement}
\begin{gather*}
\text{Cash collected}=S-\Delta AR,\qquad GP=S-\text{COGS},\qquad EBIT = GP-\text{Opex},\\
NI=(EBIT-\text{Int})(1-\tau),\qquad E_{\text{end}}=E_{\text{begin}}+NI-\text{Div}.
\end{gather*}
\[
\text{Operating leverage: } OI=(1-v)S-F,\qquad DOL=\frac{\%\Delta OI}{\%\Delta S}=1+\frac{F}{OI}.
\]
\idea Accrual principle (revenue when earned, expense when incurred) and matching principle;
depreciation spreads capex over the asset's life. COGS $\approx$ variable, other opex $\approx$ fixed.

\blk{1.7 Ratios, DuPont and ROIC}
\begin{gather*}
EPS=\frac{NI}{N},\qquad \text{margins}=\frac{GP}{S},\ \frac{OI}{S},\ \frac{NI}{S},\qquad
\text{AR days}=\frac{AR}{S/365},\qquad \text{AP days}=\frac{AP}{\text{COGS}/365},\\
\text{Coverage}=\frac{EBIT}{\text{Int}}\ \Big(\text{or }\frac{EBITDA}{\text{Int}}\Big),\qquad
\frac PE=\frac{P}{EPS},\qquad \text{Economic earnings}=NI-r_E E,
\end{gather*}
\[
ROE=\frac{NI}{E},\qquad ROA=\frac{NI}{A},\qquad
\underbrace{ROE=\frac{NI}{S}\cdot\frac{S}{A}\cdot\frac{A}{E}}_{\text{DuPont}},\qquad
ROIC=\frac{EBIT(1-\tau)}{D+E}.
\]
\[
ROE = ROIC + \frac{D}{E}\big[ROIC-r_D(1-\tau)\big],\qquad
\operatorname{sd}(ROE)=\Big(1+\frac DE\Big)\operatorname{sd}(ROIC).
\]
\idea Leverage raises ROE only if $ROIC>r_D(1-\tau)$, and always raises its risk. ROE problems:
timing (one year), risk (ignored), book values. ROIC is independent of financing; compare it with the
WACC. Diluted EPS counts options and convertibles.

\blk{1.8 Cash-flow statement and free cash flow}
\[
CFO = NI + \text{Dep} + \text{other non-cash} - \Delta AR + \Delta AP - \Delta\text{Inv},\qquad
\Delta\text{Cash}=CFO+CFI+CFF.
\]
\[
\boxed{FCF = EBIT(1-\tau)+\text{Dep}-\text{Capex}-\Delta NWC}\qquad
\boxed{FCFE = FCF - \text{Int}(1-\tau)+\Delta D}
\]
\[
CF_D=\text{Int}-\Delta D,\qquad FCFE+CF_D = FCF+\tau\,\text{Int},\qquad
FCFE = NI+\text{Dep}-\text{Capex}-\Delta NWC+\Delta D.
\]
\idea Profits $\neq$ cash. FCF (unlevered) values the enterprise (discount at WACC, or at $r_U$ in
APV); FCFE values equity (discount at $r_E$). $\Delta NWC$ counts operating items only (not cash held
for financing, not interest-bearing debt). Non-cash charges are added back, including an increase in
deferred tax liabilities (taxes expensed but not paid).

% =====================================================================
\section*{Topic 2A --- Valuation, the NPV rule and the frictionless benchmark}

\blk{Time value of money (used throughout)}
\[
PV=\frac{C_t}{(1+r)^t},\quad
\text{Perpetuity}=\frac{C}{r},\quad
\text{Growing perpetuity}=\frac{C_1}{r-g}\ (g<r),\quad
\text{Annuity}=\frac{C}{r}\Big[1-\frac{1}{(1+r)^T}\Big],
\]
\[
\text{Growing annuity}=\frac{C_1}{r-g}\Big[1-\Big(\frac{1+g}{1+r}\Big)^{T}\Big],\quad
EAR=\Big(1+\frac{APR}{k}\Big)^{k}-1,\quad
1+r_{\text{nom}}=(1+r_{\text{real}})(1+\pi).
\]

\blk{2A.1 Present value, spot rates and bonds}
\begin{gather*}
P_A=\sum_t\frac{\E[\tilde d_t]}{(1+r_t)^t},\qquad
\text{Zero: } P=\frac{F}{(1+r_T)^T}\Rightarrow r_T=\Big(\frac FP\Big)^{1/T}-1,\\
\text{Coupon bond: } P=\sum_{t=1}^{T}\frac{c}{(1+r_t)^t}+\frac{F}{(1+r_T)^T}.
\end{gather*}
\begin{gather*}
\text{YTM } y:\ P=\sum_t\frac{CF_t}{(1+y)^t},\qquad
(1+r_2)^2=(1+r_1)(1+f_{1,2}),\\
\text{HPR}=\frac{P_1-P_0+c}{P_0},\qquad
\text{Current yield}=\frac{c}{P}.
\end{gather*}
\[
\text{Duration } D_{\text{mac}}=\frac{1}{P}\sum_t t\,\frac{CF_t}{(1+y)^t},\qquad
\frac{\Delta P}{P}\approx-\frac{D_{\text{mac}}}{1+y}\,\Delta y.
\]
\idea Prices move inversely to yields; premium if $c>y$, par if $c=y$, discount if $c<y$; at a
constant yield the price pulls to par. Longer maturity and lower coupon $\Rightarrow$ more interest-rate
sensitivity. A call provision lowers the price (bond $-$ call) and raises the required yield. Credit
spread $=y_{\text{corp}}-y_{\text{Treasury}}$.

\blk{2A.2 Discount rates, CAPM, betas and comparables}
\[
r=r_f+\beta\big(\E[\tilde r_m]-r_f\big),\qquad
\beta=\frac{\operatorname{Cov}(\tilde r,\tilde r_m)}{\operatorname{Var}(\tilde r_m)},\qquad
r_A=\frac{E}{V}r_E+\frac{D}{V}r_D,\qquad
\beta_A=\frac{E}{V}\beta_E+\frac{D}{V}\beta_D,
\]
\[
r_E=r_A+\frac{D}{E}(r_A-r_D),\qquad \beta_E=\beta_A+\frac DE(\beta_A-\beta_D)\quad(\beta_D=0:\ \beta_A=\tfrac EV\beta_E).
\]
\idea Discount rate $=$ risk-free $+$ risk premium of the \emph{project}, not of the firm. A levered
comparable's $\beta_E$ overstates $\beta_A$: unlever each comparable, average $\beta_A$, then use it
(or relever at the project's target $D/E$). Expected return $\ne$ realized return.

\blk{2A.3 Which cash flows}
\[
\text{Unlevered FCF (as in 1.8)},\qquad
\text{After-tax salvage}=SV-\tau\,(SV-BV),\qquad
TV_T=\frac{FCF_{T+1}}{r-g}.
\]
\idea Incremental basis: include opportunity costs (e.g.\ value of an asset you already own) and
side effects (cannibalization); exclude sunk costs and allocated overhead that does not change.
Financing costs are \emph{not} in FCF (they enter the discount rate). Invest in NWC when sales start,
recover it at the end. Nominal flows with nominal rates, real with real.

\blk{2A.4 NPV rule and IRR}
\[
NPV_A=PV_A-c_A,\qquad 0=-c_A+\sum_t\frac{\E[\tilde d_t]}{(1+IRR)^t}.
\]
\idea Accept iff $NPV>0$. In a frictionless world NPV is the gain of existing shareholders whatever
the financing (equity, debt, internal funds). Positive NPV $=$ ``arbitrage across the firm's
boundary'' that cannot be scaled, so it survives in equilibrium. Fisher separation: investment decided
by the market rate, independent of owners' preferences. IRR pitfalls: borrowing-type flows (accept if
$IRR<r$), multiple or no IRRs (sign changes), scale and timing (mutually exclusive projects: use NPV or
the crossover rate).

\blk{2A.5 Debt vs.\ equity}
\[
\tilde C_D=\min(\tilde A,F)=F-\max(F-\tilde A,0),\qquad \tilde C_E=\max(\tilde A-F,0),\qquad
\tilde C_D+\tilde C_E=\tilde A.
\]
\idea Debt: fixed payoff, senior, maturity, covenants (technical default), no control unless in
default. Equity: residual, junior, control, limited liability. Risky debt $=$ safe bond $-$ put on the
assets; equity $=$ call on the assets with strike $F$.

\blk{2A.6--2A.7 Complete markets and state prices}
\[
\text{Price}=\sum_s p_s\,x_s,\qquad \sum_s p_s=\frac{1}{1+r_f},\qquad
\pi_s=p_s(1+r_f),\qquad \text{Price}=\frac{1}{1+r_f}\sum_s\pi_s x_s,\qquad
C-P=S-\frac{K}{1+r_f}.
\]
\idea Benchmark assumptions: competitive and complete markets, everyone borrows/lends at the same
terms, no taxes, no bankruptcy costs, financing does not change cash flows (no agency or information
problems). Law of one price; no arbitrage $\Leftrightarrow$ all $p_s>0$. Find $p_s$ by replicating
contingent claims with traded assets (solve the linear system). Value is additive: $A=D+E$ at state
prices for any split.

\blk{2A.8 Modigliani--Miller I}
\[
V_L=V_U,\qquad V(K)=\underbrace{C(K)}_{E}+\underbrace{\tfrac{K}{1+r_f}-P(K)}_{D}=S\ \text{(by put--call parity)}.
\]
\idea Proof by arbitrage: if $V_L>V_U$, sell levered claims and buy unlevered equity with
\emph{homemade leverage} (and conversely, undo leverage by buying the firm's debt). Capital structure
only splits a given cash flow.

\blk{2A.9 Modigliani--Miller II and the WACC}
\[
r_E=r_A+\frac{D}{E}(r_A-r_D),\qquad r_{WACC}=\frac EV r_E+\frac DV r_D=r_A=r_U\ \text{(no taxes)}.
\]
\idea $r_E$ rises linearly in $D/E$ while debt is safe, and more slowly once $r_D$ rises (risk moves
to creditors). A higher expected return on equity is compensation for risk, not a gain.
Conservation of value: financial transactions repackage risk; they neither add nor destroy value.

\blk{2A.10 The cost of debt is not the coupon or the yield}
\[
D=\frac{F}{1+y},\qquad D=\frac{F(1-p)+C\,p}{1+r_D}\ (p=\text{default prob.},\ C=\text{recovery})
\ \Rightarrow\ y>r_D\ \text{if } p>0.
\]
\idea The yield is the \emph{promised} return; the expected return $r_D$ is lower when default is
possible. Using $y$ as $r_D$ overstates the cost of debt.

% =====================================================================
\section*{Topic 2B --- Taxes, bankruptcy costs and the static trade-off}

\blk{2B.1--2B.2 Interest tax shield and MM I with taxes}
\[
ITS_t=\tau\,r_D D_{t-1}=\tau\times\text{Interest},\qquad
\text{payout to investors}=(1-\tau)X+\tau r_D D,\qquad
\boxed{V_L=V_U+PV(ITS)}.
\]
\[
\text{Market-value balance sheet: } V_U+V_{ITS}=V_L=D+E,\qquad
\frac DV r_D+\frac EV r_E=\frac{V_U}{V_L}r_U+\frac{V_{ITS}}{V_L}r_{ITS}.
\]
\idea The friction is the differential treatment of interest (deductible) vs.\ dividends. Unlevered
FCF ignores the shield; it enters through the WACC or the APV. The gain accrues to shareholders.
Debt-financed repurchase: on announcement $P=\dfrac{V_U+PV(ITS)}{N_0}$; shares bought $=D/P$.

\blk{PV of the tax shield and MM II with taxes, by debt policy}
\begin{center}\small
\begin{tabular}{lll}
\hline
Debt policy & $V_{ITS}$ & $r_E$ \\ \hline
Constant (perpetual) debt, $r_{ITS}=r_D$ & $\tau D$ & $r_U+\frac DE(1-\tau)(r_U-r_D)$;\ \ $\beta_E=\beta_U[1+(1-\tau)\frac DE]$ \\
Continuous target $D/V$, $r_{ITS}=r_U$ & $\dfrac{\tau r_D D}{r_U-g}$ & $r_U+\frac DE(r_U-r_D)$ \\
Annual target $D/V$ & $\dfrac{\tau r_D D}{r_U-g}\cdot\dfrac{1+r_U}{1+r_D}$ & $r_U+\frac DE(r_U-r_D)$ (approx.) \\ \hline
\end{tabular}
\end{center}

\blk{2B.3--2B.4 After-tax WACC and APV}
\begin{gather*}
\boxed{r_{WACC}=\frac EV r_E+\frac DV r_D(1-\tau)}\\
r_{WACC}=r_U-\tau r_D\frac DV\ \text{(continuous target)},\qquad
r_U-\tau r_D\frac DV\frac{1+r_U}{1+r_D}\ \text{(annual target)},\\
r_U\Big(1-\tau\frac DV\Big)\ \text{(constant debt, } g=0).
\end{gather*}
\[
r_{WACC}-g=(r_U-g)\Big(1-\frac{V_{ITS}}{V_L}\Big),\qquad
APV:\ V_L=V_U+PV(ITS)-PV(\text{distress costs})+PV(\text{other side effects}).
\]
\idea WACC method: unlever comparables to get $r_U$, relever at the target $D/E$, compute $r_E$ and the
WACC, discount FCF. It assumes a constant $D/E$. APV: discount FCF at $r_U$, add the shield at its
own risk; best when the debt schedule is known (e.g.\ a fixed loan). Both agree under constant $D/E$.
Risky debt: the shield uses the coupon actually paid, $\bar r_D=\E[C_{t+1}]/D_t$; after-tax cost of
debt $=r_D-\tau\bar r_D$.

\blk{2B.5 The 100\%-debt puzzle and personal taxes}
\[
\frac{V_L}{V_U}=1+\frac{\tau D}{V_U},\qquad
\text{Miller: } G_L=\Big[1-\frac{(1-\tau_c)(1-\tau_s)}{1-\tau_b}\Big]D
\]
($\tau_c$ corporate, $\tau_s$ and $\tau_b$ personal tax rates on equity and on debt income).
\idea MM with corporate taxes alone implies all-debt firms; observed leverage is moderate
($D/(D+E)\approx$ 0.1--0.45 by industry). Personal taxes on interest reduce, and can cancel, the gain.

\blk{2B.6 Bankruptcy costs}
\idea Default is a transfer of control to creditors, not an MM violation. Costs: \emph{direct}
(legal, administrative, management time, fire sales) and \emph{indirect} (lost customers, suppliers
cutting credit, workers leaving, agency conflicts: asset substitution, debt overhang, cashing out).
Expected costs are borne by shareholders ex ante through a lower price of the claims they issue.

\blk{2B.7 Static trade-off model}
\begin{gather*}
V_U=\frac{(1-\tau)\E[X]}{1+r_f},\qquad
V(K)=V_U+\frac{1}{1+r_f}\Big[\tau\Big(\int_{\underline X}^{K}x\,h(x)\,dx+K\big(1-H(K)\big)\Big)-C\,H(K)\Big],\\
\text{FOC: }\ \tau\big[1-H(K^*)\big]=C\,h(K^*),\qquad
\frac{dK^*}{d\tau}>0,\ \ \frac{dK^*}{dC}<0,\\
\text{uniform } h:\ K^*=\bar X-\frac{C}{\tau}\ \text{ if } C<\tau(\bar X-\underline X),\ \text{else } K^*=\underline X.
\end{gather*}
\idea Safe debt ($K<\underline X$): $V=V_U+\tau K/(1+r_f)$. Marginal benefit: shield on the solvent
states; marginal cost: expected bankruptcy cost at the margin. Same logic as
$V_L=V_U+PV(ITS)-PV(\text{bankruptcy costs})$.

\blk{2B.8 Bank capital debate (Admati et al.)}
\idea ``More equity lowers ROE'': expected ROE falls, but so does its risk; required $r_E$ falls (MM II)
and the WACC is unchanged. ``Equity is expensive'': only privately, through the lost tax shield, which
is a transfer from taxpayers, not a social cost. Bank failure has a social cost (externality), so
the social optimum has more equity than banks choose.

\blk{2B.9 Payout irrelevance (MM3)}
\[
\text{Sources}=\text{Uses}:\ X_t+\text{New issues}_t=I_t+\text{Div}_t,\qquad
V_t=\E\Big[\sum_{i\ge1}\delta^{i}\big(X_{t+i}-I_{t+i}\big)\Big].
\]
\idea With the investment plan fixed, value does not depend on payout: a higher dividend is financed
by issuing claims of equal value. Dividend vs.\ repurchase: same total value to shareholders in a
perfect market (ex-dividend price drops by the dividend).

% =====================================================================
\section*{Exam traps}
\begin{itemize}
\item Negative book equity is not insolvency; profit is not cash; depreciation is not a source of cash
      (only the tax it saves is).
\item $EV=E^{\text{mkt}}+D-\text{Cash}$, not market cap plus total liabilities. Always say which
      $D/E$ (book or market) you use.
\item FCF taxes EBIT as if the firm were unlevered: interest and the tax shield stay out of FCF (the
      shield enters the WACC or the APV).
\item $\Delta NWC$ counts operating items only, not excess cash and not short-term debt.
\item A higher ROE is not better performance: compare $ROIC$ with $r_D(1-\tau)$ and account for the
      added risk.
\item Use comparables' \emph{asset} betas; never discount a project at the firm's $r_E$ or WACC
      unless the project has the firm's risk and financing.
\item ``Debt is cheaper'' is wrong without taxes: $r_E$ rises to offset it (MM II) and the WACC is
      unchanged.
\item Equity gets riskier with leverage even when debt is riskless: the same risk falls on fewer
      claims.
\item The yield (or the coupon) is not $r_D$ when default is possible.
\item State prices are not probabilities: $\sum_s p_s=1/(1+r_f)$, and $\pi_s$ are risk-neutral, not
      real-world, probabilities.
\item IRR does not rank projects (scale, timing, sign changes): use NPV.
\item $V_{ITS}=\tau D$ only for permanent, constant, safe debt; with a target $D/E$ discount the
      shield at $r_U$.
\item $r_{WACC}=r_U(1-\tau D/V)$ only for constant debt and $g=0$.
\item Bankruptcy \emph{costs}, not default itself, break MM; shareholders bear them ex ante.
\end{itemize}

\end{document}
